% from ega3-1-fr.tex
\oldpage[0\textsubscript{III}]{33}
\begin{enumerate}
  \item[(i)] Pour chaque indice $i$, les suites
    \[
    \begin{aligned}
      0\longrightarrow{}&K^i\xrightarrow{\varepsilon}L^{i0}
        \longrightarrow L^{i1}\longrightarrow\cdots\\
      0\longrightarrow{}&\mathrm B^i(K^\bullet)\xrightarrow{\varepsilon}
        \mathrm B_{\mathrm I}^i(L^{\bullet,0})
        \longrightarrow\mathrm B_{\mathrm I}^i(L^{\bullet,1})
        \longrightarrow\cdots\\
      0\longrightarrow{}&\mathrm Z^i(K^\bullet)\xrightarrow{\varepsilon}
        \mathrm Z_{\mathrm I}^i(L^{\bullet,0})
        \longrightarrow\mathrm Z_{\mathrm I}^i(L^{\bullet,1})
        \longrightarrow\cdots\\
      0\longrightarrow{}&\mathrm H^i(K^\bullet)\xrightarrow{\varepsilon}
        \mathrm H_{\mathrm I}^i(L^{\bullet,0})
        \longrightarrow\mathrm H_{\mathrm I}^i(L^{\bullet,1})
        \longrightarrow\cdots
    \end{aligned}
    \]
    sont exactes, autrement dit $(L^{i,\bullet})$,
    $(\mathrm B_{\mathrm I}^i(L^{\bullet,\bullet}))$,
    $(\mathrm Z_{\mathrm I}^i(L^{\bullet,\bullet}))$ et
    $(\mathrm H_{\mathrm I}^i(L^{\bullet,\bullet}))$ sont respectivement
    des \emph{résolutions} de $K^i$, $\mathrm B^i(K^\bullet)$,
    $\mathrm Z^i(K^\bullet)$ et $\mathrm H^i(K^\bullet)$.
  \item[(ii)] Pour chaque $j$, le complexe simple $L^{\bullet,j}$ est
    \emph{scindé}, autrement dit, les suites exactes
    \[
      \label{0.11.4.2.1-fr}
      0\longrightarrow\mathrm B_{\mathrm I}^i(L^{\bullet,j})
      \longrightarrow\mathrm Z_{\mathrm I}^i(L^{\bullet,j})
      \longrightarrow\mathrm H_{\mathrm I}^i(L^{\bullet,j})
      \longrightarrow0
      \tag{11.4.2.1}
    \]
    \[
      \label{0.11.4.2.2-fr}
      0\longrightarrow\mathrm Z_{\mathrm I}^i(L^{\bullet,j})
      \longrightarrow L^{ij}
      \longrightarrow\mathrm B_{\mathrm I}^{i+1}(L^{\bullet,j})
      \longrightarrow0
      \tag{11.4.2.2}
    \]
    sont scindées.
\end{enumerate}

On prouve (M, XVII, 1.2) que si tout objet de $\mathcal C$ est sous-objet
d'un objet injectif, tout complexe $K^\bullet$ de $\mathcal C$ admet une
résolution de Cartan-Eilenberg \emph{injective}, c'est-à-dire formée d'objets
injectifs $L^{ij}$ (la condition (ii) ci-dessus entraîne alors que les
$\mathrm B_{\mathrm I}^i(L^{\bullet,j})$,
$\mathrm Z_{\mathrm I}^i(L^{\bullet,j})$ et
$\mathrm H_{\mathrm I}^i(L^{\bullet,j})$ sont aussi des objets injectifs).
En outre, pour tout morphisme $f:K^\bullet\to K'^\bullet$ de complexes de
$\mathcal C$, toute résolution de Cartan-Eilenberg $L^{\bullet\bullet}$ de
$K^\bullet$ et toute résolution injective de Cartan-Eilenberg
$L'^{\bullet\bullet}$ de $K'^\bullet$, il existe un morphisme de bicomplexes
$F:L^{\bullet\bullet}\to L'^{\bullet\bullet}$ compatible avec $f$ et les
augmentations, et si $f$ et $g$ sont deux morphismes homotopes de $K^\bullet$
dans $K'^\bullet$, les morphismes correspondants de $L^{\bullet\bullet}$
dans $L'^{\bullet\bullet}$ sont homotopes (\emph{loc. cit.}).

Lorsque $K^\bullet$ est \emph{limité inférieurement} (resp.
\emph{supérieurement}), on peut prendre $L^{\bullet\bullet}$ tel que
$L^{ij}=0$ pour $i<i_0$ (resp. $i>i_0$) si $K^i=0$ pour $i<i_0$ (resp.
$i>i_0$) (M, XVII, 1.3).

Supposons d'autre part qu'il existe un entier $n$ tel que tout objet de
$\mathcal C$ admette une \emph{résolution injective de longueur $\leq n$}~;
alors on peut supposer que l'on a $L^{ij}=0$ pour $j>n$ (M, XVII, 1.4).
\end{env}

\begin{env}[11.4.3]
\label{0.11.4.3-fr}
Soit maintenant $T$ un \emph{foncteur covariant additif} de $\mathcal C$
dans une catégorie abélienne $\mathcal C'$. Étant donnés un complexe
$K^\bullet$ de $\mathcal C$ et une résolution de Cartan-Eilenberg injective
$L^{\bullet\bullet}$ de $K^\bullet$, supposons que le complexe (simple) défini
par le bicomplexe $T(L^{\bullet\bullet})$ existe
(cf. \hyperref[0.11.3.1-fr]{11.3.1})~; alors les deux suites spectrales
$'E(T(L^{\bullet\bullet}))$ et $''E(T(L^{\bullet\bullet}))$ de ce bicomplexe
sont dites \emph{suites spectrales d'hypercohomologie de $T$ par rapport au
complexe $K^\bullet$}~; en vertu de
(\hyperref[0.11.4.2-fr]{11.4.2}) et
(\hyperref[0.11.3.2-fr]{11.3.2}), elles ne dépendent effectivement que de
$K^\bullet$ et non de la résolution de Cartan-Eilenberg injective
$L^{\bullet\bullet}$ choisie~; en outre, elles dépendent
\emph{fonctoriellement} de $K^\bullet$. Elles ont un même aboutissement
$\mathrm H^\bullet(T(L^{\bullet\bullet}))$ appelé encore
l'\emph{hypercohomologie de $T$ par rapport à $K^\bullet$}, et noté
$\mathrm R^\bullet T(K^\bullet)$. On montre que les termes $E_2$ des deux
suites spectrales précédentes sont donnés par
\[
  \label{0.11.4.3.1-fr}
  'E_2^{pq}=\mathrm H^p(\mathrm R^qT(K^\bullet))
  \tag{11.4.3.1}
\]
\[
  \label{0.11.4.3.2-fr}
  ''E_2^{pq}=\mathrm R^pT(\mathrm H^q(K^\bullet))
  \tag{11.4.3.2}
\]
